\chapter{CONCLUSION AND FUTURE WORK}
\label{chap:conclusion}

This chapter synthesizes the theoretical, methodological, and empirical contributions consolidated throughout this doctoral research. Grounded in the Design Science Research Methodology (DSRM) \cite{peffers_2006, hevner_2004} and classical Control Theory duality \cite{kalman_1959}, this study addressed the problem of observability fragmentation and root-cause ambiguity in Kubernetes-orchestrated cloud-native microservice environments. We formally answer the formulated Research Questions, summarize the empirical validation of our research hypotheses ($H_{1,1}$, $H_{1,2}$, and $H_{1,3}$), consolidate the architectural knowledge crystallized in our Control Principles, Contingency Heuristics, and Markdown Architectural Decision Records (\textbf{ADR-0001}--\textbf{ADR-0006}), and delineate the roadmap of extensions leading up to the final thesis defense.

\section{Synthesis of Research Questions and Hypothesis Validation}
\label{sec:conclusion-synthesis}

The primary objective of this research---to propose, instantiate, and empirically validate a vendor-agnostic, open-source observability reference architecture for Kubernetes (\textbf{PANOPTES})---was fully achieved. Table~\ref{tab:rq_hypothesis_synthesis} maps each Secondary Research Question (\textbf{SRQ1}--\textbf{SRQ4}), Specific Objective (\textbf{SO1}--\textbf{SO4}), and Research Hypothesis ($H_{1,1}$--$H_{1,3}$) to its corresponding empirical outcome.

\begin{table}[htbp]
    \centering
    \caption{Synthesis of Research Questions, Hypotheses, and Empirical Validation Outcomes.}
    \label{tab:rq_hypothesis_synthesis}
    \footnotesize
    \renewcommand{\arraystretch}{1.2}
    \begin{tabularx}{\textwidth}{@{} l p{2.8cm} X p{3.6cm} @{}}
        \toprule
        \textbf{Construct} & \textbf{Focus / Metric} & \textbf{Key Empirical / Architectural Finding} & \textbf{Validation Status \& Evidence} \\
        \midrule
        \textbf{SRQ1 / SO1} & Minimum Architectural Requirements & Formalized the four-tier observability pipeline (Agent, Aggregation/Routing, Specialized Datastore, and Query/Correlation tiers) over the telemetry tuple $O = \{M, L, T, E\}$ via OpenTelemetry, Prometheus, Loki, and Tempo. & \textbf{Validated} (Chapters~\ref{chap:related-work}, \ref{chap:cloud-computing}, and~\ref{chap:experiment-implementation}) \\
        \addlinespace
        \textbf{SRQ2 / SO2} & Semantic Integration vs.\ Resource Footprint & Enforced unified OTLP/Kubernetes resource labels (\texttt{CP2}), metric exemplars, trace-to-log bindings (\texttt{CP1}), label-only log indexing, and object-block trace storage under declarative IaC (\texttt{ConH1}--\texttt{ConH3}). & \textbf{Validated} (Chapters~\ref{chap:methodology} and~\ref{chap:experiment-implementation}; \textbf{ADR-0001}--\textbf{ADR-0006}) \\
        \addlinespace
        \textbf{SRQ3 / SO3} & Diagnostic Effectiveness under Chaos & Evaluated across $N=30$ paired, seeded Chaos Mesh trials (\texttt{cpu-stress}, \texttt{pod-kill}, \texttt{network-delay}) comparing Vanilla CLI (S3/S3-Strat), Monomodal Tracing (S4a), and PANOPTES Multimodal RCA (S4b). & \textbf{Validated} (Chapter~\ref{chap:results-and-discussion}; Tables~\ref{tab:s3_s4_comparison} and~\ref{tab:stratified_breakdown}) \\
        \addlinespace
        \textbf{SRQ4 / SO4} & Evolutionary Stages \& Maturity Gaps & Synthesized $N=10$ primary studies (from $312$ initial records) via PRISMA Rapid Review and formalized the 5-level Observability Maturity Model (OMM), empirically validating Level~4 against Levels~2 and~3. & \textbf{Validated} (Chapters~\ref{chap:related-work} and~\ref{chap:results-and-discussion}) \\
        \midrule
        \textbf{Hypothesis $H_{1,1}$} & Diagnostic Efficiency ($MTTDiag$) & PANOPTES reduced median $MTTDiag$ by $\mathbf{94.24\%}$ (from $300.00\,\text{s}$ in Vanilla S3 to $\mathbf{17.27\,\text{s}}$ in S4b; mean reduced by $77.51\%$ from $242.42\,\text{s}$ to $54.52\,\text{s}$). & \textbf{Accepted ($H_{0,1}$ Rejected):} Mann-Whitney $p = 4.55 \times 10^{-6}$; Wilcoxon $p = 3.06 \times 10^{-5}$; Log-Rank $p = 1.34 \times 10^{-6}$; $\hat{A}_{12} = 0.817$ (Large Effect). \\
        \addlinespace
        \textbf{Hypothesis $H_{1,2}$} & Diagnostic Accuracy ($Acc$) & PANOPTES elevated root-cause accuracy from $20.00\%$ ($6/30$ in S3), $23.33\%$ ($7/30$ in S3-Strat), and $0.00\%$ ($0/30$ in S4a) to $\mathbf{83.33\%}$ ($\mathbf{25/30}$ in S4b; $100\%$ on \texttt{cpu-stress} and $100\%$ on \texttt{pod-kill}). & \textbf{Accepted ($H_{0,2}$ Rejected):} Fisher's Exact $\text{OR} = 20.0, \; p = 7.74 \times 10^{-7}$; Paired McNemar $p = 7.83 \times 10^{-5}$. \\
        \addlinespace
        \textbf{Hypothesis $H_{1,3}$} & Bounded Resource Overhead ($\Delta_{\text{CPU}}, \Delta_{\text{RAM}}$) & Steady-state PANOPTES CPU consumed $\mathbf{9.51\% \pm 2.14\%}$ ($189.47\,\text{mCPU}$, med.\ $8.83\%$) of workload CPU ($\le 15\%$), while backend RSS memory ($3{,}235.73\,\text{MiB}$) occupied $\mathbf{26.33\%}$ of the $12\,\text{GiB}$ cluster capacity ($\le 30\%$). & \textbf{Accepted ($H_{0,3}$ Rejected):} One-Sample Wilcoxon $p = 6.52 \times 10^{-9}$ (CPU) and $p = 9.31 \times 10^{-10}$ (RAM). \\
        \bottomrule
    \end{tabularx}
    \par\medskip\ABNTEXfontereduzida\selectfont\textbf{Source:} The author (2026) \par\medskip
\end{table}


Based on the evidence consolidated in Chapters~\ref{chap:related-work} through~\ref{chap:results-and-discussion}, the research questions are answered as follows:

\begin{itemize}
    \item \textbf{Answer to SRQ1 (Minimum Architectural Requirements):} Through the Rapid Review of the state of the art (Chapter~\ref{chap:related-work}) and technical analysis of the CNCF ecosystem (Chapter~\ref{chap:cloud-computing}), we established that a complete Kubernetes observability layer must capture the four-dimensional telemetry tuple $O = \{M, L, T, E\}$ across four functional tiers: (1)~a non-intrusive \textit{Agent Collection Tier} (OpenTelemetry Collector DaemonSets and kubelet stats receivers); (2)~an \textit{Aggregation and Normalization Tier} enforcing unified OTLP resource attributes; (3)~a \textit{Specialized Low-Overhead Datastore Tier} pairing pull-based time-series scraping (Prometheus), label-only chunked log storage (Grafana Loki), and index-free object trace blocks (Grafana Tempo); and (4)~a \textit{Cross-Pillar Correlation and RCA Tier} exposing unified programmatic APIs (PromQL, LogQL, and TraceQL) alongside visual dashboards (Grafana).

    \item \textbf{Answer to SRQ2 (Semantic Integration within a Bounded Resource Envelope):} Semantic integration without excessive cluster burden is achieved by replacing heavy full-text log inverted indices and secondary trace indices with shared metadata bindings (\texttt{CP1} and \texttt{CP2}): Prometheus histograms bind to Tempo traces via OTLP exemplars, Tempo spans pivot to Loki streams via contextual \texttt{trace\_id} and pod labels, and collector pipelines apply batching (\texttt{ConH3}). In our steady-state evaluation across $N=30$ anti-pseudoreplicated samples ($16\,\text{s}$ interval, \textbf{ADR-0003}), the PANOPTES stack consumed only $\mathbf{9.51\% \pm 2.14\%}$ (median $\mathbf{8.83\%}$) of the active workload's CPU utilization ($189.47\,\text{mCPU}$ vs.\ $1{,}994.43\,\text{mCPU}$, $p = 6.52 \times 10^{-9}$ against $\theta_{\text{CPU}} = 15\%$), while its centralized storage backends maintained a stable working set of $3{,}235.73\,\text{MiB}$ ($\mathbf{26.33\%}$ of the $12\,\text{GiB}$ cluster RAM capacity, $p = 9.31 \times 10^{-10}$ against $\theta_{\text{RAM}} = 30\%$), confirming \textbf{Hypothesis $H_{1,3}$}.

    \item \textbf{Answer to SRQ3 and Main Research Question (Diagnostic Effectiveness vs.\ Kubernetes Baseline):} Across $N=30$ paired, stratified Chaos Mesh fault-injection trials (\texttt{RANDOM\_SEED = 42}, \textbf{ADR-0001}), native Kubernetes CLI inspection (\textbf{Scenario~3} and its independent replication \textbf{Scenario~3-Stratified}) suffered $73.33\%\text{--}80.00\%$ right-censoring at $300\,\text{s}$, failing completely on \texttt{cpu-stress} ($0/11$) and achieving only $20.00\%\text{--}23.33\%$ global accuracy. Crucially, evaluating distributed tracing in isolation without self-time decomposition (\textbf{Scenario~4a}) yielded $\mathbf{0.00\%}$ accuracy ($30/30$ misattributions) due to upstream cascading span duration bias. By contrast, the **PANOPTES Multimodal Correlated RCA** (\textbf{Scenario~4b})---combining PromQL saturation/reset detection, LogQL exception streams, and Tempo \textit{exclusive span self-time decomposition} ($T_{\text{self}} = T_{\text{span}} - \sum T_{\text{children}}$) with L4/L7 peer attribution (\textbf{ADR-0002})---elevated diagnostic accuracy to $\mathbf{83.33\%}$ ($25/30$; $100\%$ on \texttt{cpu-stress} and $100\%$ on \texttt{pod-kill}; Fisher $p = 7.74 \times 10^{-7}$, McNemar $p = 7.83 \times 10^{-5}$, confirming \textbf{Hypothesis $H_{1,2}$}) and reduced median $MTTDiag$ by $\mathbf{94.24\%}$ from $300.00\,\text{s}$ to $\mathbf{17.27\,\text{s}}$ (Mann-Whitney $p = 4.55 \times 10^{-6}$, Wilcoxon $p = 3.06 \times 10^{-5}$, Log-Rank $p = 1.34 \times 10^{-6}$, Vargha-Delaney $\hat{A}_{12} = 0.8167$, confirming \textbf{Hypothesis $H_{1,1}$}).

    \item \textbf{Answer to SRQ4 (Evolutionary Classification and Maturity Validation):} We formalized the 5-stage \textbf{Observability Maturity Model (OMM)} in Chapter~\ref{chap:related-work} following the methodological guidelines of \citeonline{becker_2009} and \citeonline{debruin_2005}. Directly addressing the ``Validation Gap'' identified by \citeonline{wendler_2012}, our factorial experiment empirically demonstrated the operational progression across OMM stages: proving why \textbf{Level~2 (Initial/Ad-Hoc CLI)} and \textbf{Level~3 (Managed Silos / Monomodal Tracing)} fail under distributed microservice faults, and establishing \textbf{Level~4 (Optimized / Correlated PANOPTES)} as the mandatory diagnostic substrate required to support future \textbf{Level~5 (Predictive/Autonomic AIOps)} control loops.
\end{itemize}

\section{Consolidated Architectural and Methodological Knowledge (DSR Knowledge Base)}
\label{sec:consolidated-decisions}

Beyond the concrete Infrastructure as Code instantiation, a core contribution of Design Science Research is the codification of reusable design knowledge \cite{gregor_2013}. Through the iterative Design Cycle and our six Markdown Architectural Decision Records (\textbf{ADR-0001}--\textbf{ADR-0006}, Section~\ref{sec:madr-decisions}), this research consolidated five foundational lessons for cloud-native systems engineering:

\begin{enumerate}
    \item \textbf{Multimodal Self-Time Decomposition is Essential to Overcome Cascading Span Bias (\texttt{CP1}, \texttt{CP2}, ADR-0002):} In synchronous microservice call graphs, total span duration ($T_{\text{span}}$) monotonically increases toward the ingress gateway. Accurate automated RCA requires subtracting child span intervals ($T_{\text{self}} = T_{\text{span}} - \sum T_{\text{children}}$), resolving L4/L7 client-to-peer attributes (\texttt{server.address}, \texttt{net.peer.name}), and cross-referencing infrastructure saturation (PromQL) and process termination signals (LogQL).
    \item \textbf{Separation of Proportional CPU Overhead from Fixed Storage Working Sets (\texttt{ConH3}, ADR-0005):} While node-level telemetry collection scales lightweightly with workload CPU ($\sim 9.5\%$), in-cluster TSDB and trace backends require a fixed baseline RAM allocation ($\sim 3.2\,\text{GiB}$). In memory-constrained clusters ($4\,\text{GiB}$/node), pruning non-transactional bloat workloads (\textbf{ADR-0005}) or segregating storage backends onto dedicated infrastructure nodes is critical to prevent kernel swap thrashing (\texttt{kswapd0}).
    \item \textbf{Deterministic Virtualization Parity via K3d and Linux \texttt{cgroups} (ADR-0006):} Migrating across Vagrant/VirtualBox (Type-2 NEM I/O degradation) and Multipass/Hyper-V (cross-VM \texttt{netsh portproxy} tunnel drops) to containerized multi-node K3s (\textbf{K3d}) on native WSL2 with hard per-node \texttt{cgroup} limits (\texttt{--cpus="2.0" --memory="4g" --memory-swap="4g"}) preserved strict 3-node hardware isolation while guaranteeing sub-millisecond API polling stability across 120+ hours of Chaos Mesh testing.
    \item \textbf{Anti-Pseudoreplication and Paired Block Designs in Cloud Benchmarks (ADR-0001, ADR-0003, ADR-0004):} Aligning resource polling ($16\,\text{s}$) strictly above the Kubernetes \texttt{metrics-server} scraping window ($15\,\text{s}$) eliminates step-function autocorrelation \cite{hurlbert_1984}, while pod-instance chaos binding and seeded block schedules (\texttt{RANDOM\_SEED = 42}) enable high-power paired non-parametric inference (Wilcoxon and McNemar tests).
\end{enumerate}

\section{Limitations, Future Work, and Post-Qualification Roadmap}
\label{sec:future-work}

While the controlled laboratory evaluation on the 3-node K3d cluster successfully guaranteed internal, construct, and statistical conclusion validity, several exciting avenues remain to expand external validity and advance the architecture toward OMM Level~5 leading up to the final thesis defense:

\begin{enumerate}
    \item \textbf{External Validity Extension on Public Cloud Infrastructure (GCP/GKE via Terraform):} Because the PANOPTES IaC architecture strictly decouples node provisioning from Kubernetes software configuration (Ansible, Helm, and Python RCA harnesses), the next experimental phase will provision a multi-node managed cluster on \textbf{Google Cloud Platform (Google Kubernetes Engine --- GKE)} using \textbf{Terraform}. Replicating the $2\times2$ factorial experiment on GKE will allow us to quantify the impact of real cloud CNI fabrics, cloud block storage IOPS, and multi-tenant infrastructure variance against our deterministic K3d baseline.
    
    \item \textbf{Refinement of Sub-Threshold Asynchronous Latency Detection:} In Scenario~4b, the $4$ right-censored trials occurred exclusively during \texttt{network-delay} perturbations on lower-frequency checkout paths (\texttt{payment} and \texttt{shipping}) when synthetic user traffic generated fewer post-$T_0$ spans than required by the conservative multi-vote threshold. Future work will integrate Prometheus histogram exemplar triggers ($p95/p99$ rate-of-change alerts) directly into the Tempo self-time voter to accelerate convergence on low-throughput RPC edges.
    
    \item \textbf{Kernel-Level Zero-Code Instrumentation via eBPF (Cilium / Grafana Beyla):} To operationalize Pillar~1 of our Strategic Research Roadmap (Figure~\ref{fig:research-roadmap}), we plan to integrate eBPF kernel probes to capture socket-level TCP retransmissions and round-trip latency directly from the Linux kernel without requiring application-level OpenTelemetry SDKs.
    
    \item \textbf{Closing the Control Loop toward OMM Level~5 (Autonomic Self-Healing):} Currently, PANOPTES fulfills the \textit{Sensor} and \textit{Diagnostic Controller} stages of Kalman's feedback loop (OMM Level~4). A natural extension is coupling the deterministic output of the PANOPTES RCA engine (\texttt{culprit}, \texttt{detection\_modality}, and \texttt{evidence\_id}) to a Kubernetes Operator actuator capable of executing automated remediation (e.g., rolling pod restarts, circuit breaking, or horizontal autoscaling) and dynamic tail-based sampling (\texttt{ConH3}).
\end{enumerate}

\section{Schedule and Milestones Toward Final Defense}
\label{sec:schedule-milestones}

With the core instantiation of PANOPTES, the Observability Maturity Model (OMM), and the full local $2\times2$ factorial evaluation ($N=30$ across Scenarios~1, 2, 3, 3-Stratified, 4a, and~4b) completed and statistically validated in this qualification document, the remaining schedule leading up to the final doctoral defense is structured as follows:

\begin{itemize}
    \item \textbf{Phase 1 (Months 1--2 Post-Qualification):} Incorporation of the Qualification Committee's recommendations and execution of the complementary external-validity campaign on Google Cloud Platform (GKE provisioned via Terraform).
    \item \textbf{Phase 2 (Month 3):} Integration of the cross-environment comparative analysis (Local K3d \texttt{cgroup} testbed vs.\ Public Cloud GKE testbed) and refinement of the Tempo exemplar-assisted latency heuristic.
    \item \textbf{Phase 3 (Month 4):} Submission of the consolidated empirical findings (PANOPTES Architecture, OMM validation, and Multimodal Self-Time RCA benchmark) to a high-impact peer-reviewed journal (e.g., \textit{Journal of Systems and Software}, \textit{IEEE Transactions on Services Computing}, or \textit{Empirical Software Engineering}).
    \item \textbf{Phase 4 (Months 5--6):} Final manuscript revision, open-source replication package publication, and Final Doctoral Thesis Defense.
\end{itemize}
